\documentclass[11pt,a4paper]{article}
\usepackage{fontspec}
\usepackage{geometry}
\usepackage{amsmath,amssymb,mathtools}
\usepackage{booktabs}
\usepackage{enumitem}
\usepackage{xcolor}
\usepackage{titlesec}
\usepackage{fancyhdr}
\usepackage[unicode=true]{hyperref}
\usepackage[most]{tcolorbox}
\usepackage{lastpage}
\usepackage{microtype}
\usepackage{xeCJK}
\geometry{top=22mm,bottom=24mm,left=24mm,right=24mm,headheight=14pt}
\definecolor{ink}{HTML}{172033}\definecolor{accent}{HTML}{315A8A}\definecolor{accentlight}{HTML}{EAF1F8}\definecolor{rulegray}{HTML}{D7DEE8}\definecolor{muted}{HTML}{5C6678}
\setmainfont{DejaVu Serif}\setsansfont{DejaVu Sans}\setmonofont{DejaVu Sans Mono}\setCJKmainfont{Noto Serif CJK JP}\setCJKsansfont{Noto Sans CJK JP}
\hypersetup{colorlinks=true,linkcolor=accent,urlcolor=accent,citecolor=accent,pdftitle={法 3 数列における半周期対称性},pdfsubject={数学的形式化とアルゴリズム監査},pdfkeywords={漸化式, 法3, 反周期性, フーリエ, コンパニオン行列}}
\pagestyle{fancy}\fancyhf{}\fancyhead[L]{\small\color{muted}法 3 の半周期対称性}\fancyhead[R]{\small\color{muted}形式化と監査}\fancyfoot[C]{\small\color{muted}\thepage\,/\,\pageref*{LastPage}}\renewcommand{\headrulewidth}{0.3pt}\renewcommand{\footrulewidth}{0pt}
\titleformat{\section}{\Large\bfseries\color{ink}}{\thesection.}{0.6em}{}[\vspace{0.2em}\color{rulegray}\titlerule]\titleformat{\subsection}{\large\bfseries\color{accent}}{\thesubsection}{0.6em}{}\titlespacing*{\section}{0pt}{2.2ex plus .5ex minus .2ex}{1.3ex}\titlespacing*{\subsection}{0pt}{1.7ex plus .4ex minus .2ex}{0.8ex}
\setlist[itemize]{leftmargin=1.6em,itemsep=0.25em,topsep=0.4em}\setlist[enumerate]{leftmargin=1.8em,itemsep=0.25em,topsep=0.4em}\setlength{\parindent}{0pt}\setlength{\parskip}{0.55em}\allowdisplaybreaks
\newtcolorbox{summarybox}{colback=accentlight,colframe=accent,boxrule=0.6pt,arc=2mm,left=4mm,right=4mm,top=3mm,bottom=3mm,before skip=1em,after skip=1em}\newcommand{\F}{\mathbb F}\newcommand{\Z}{\mathbb Z}\newcommand{\concat}{\mathbin{\Vert}}
\begin{document}
\begin{titlepage}\thispagestyle{empty}\vspace*{22mm}{\color{accent}\rule{\textwidth}{1.2pt}}\\[15mm]{\fontsize{27}{33}\selectfont\bfseries\color{ink}法 3 数列における半周期対称性\par}\vspace{9mm}{\Large\color{accent}数学的形式化とアルゴリズム監査\par}\vspace{13mm}\begin{summarybox}本稿は、しばしば混同される三つの現象、すなわちサイクルの符号反転、半周期反周期性、反転を厳密に区別する。非零原始語に対する唯一の反対シフトを証明し、組合せ的・スペクトル的・行列的特徴付けを与え、これらをモジュラー漸化式の全探索解析器に結び付ける。\end{summarybox}\vfill{\large\color{muted}リポジトリ: \texttt{Modular-recurrence-symetry-analyzer-}\par}{\large\color{muted}統合 PDF 版 - 2026 年 8 月 5 日\par}\vspace{12mm}{\color{accent}\rule{\textwidth}{1.2pt}}\end{titlepage}
\renewcommand{\contentsname}{目次}\tableofcontents\clearpage
\section{目的}
本稿は、リポジトリ \href{https://github.com/59200/k-bonacci-pisano-finder}{\texttt{59200/k-bonacci-pisano-finder}} で「3 に対して相補的な周期」「二つに切ると自己相補的」と非形式的に呼ばれている現象を定式化する。

不変な枠組みは体 \(\F_3=\{0,1,2\}\) 上の周期語であり、加法的対合は \(\nu(x)=-x\pmod3\)、すなわち \(\nu(0)=0\)、\(\nu(1)=2\)、\(\nu(2)=1\) である。したがって「3 の補数」は \textbf{法 3 における加法逆元} または \textbf{法 3 の符号反転による補元} と呼ぶべきである。0 は 0 に写るため、整数代表の十進和が常に 3 になるわけではない。
\section{三つの異なる性質}
\(w=(w_0,\ldots,w_{T-1})\in\F_3^T\) を原始周期 \(T\) の語とする。
\subsection{反対周期の対}
すべての \(n\) について \(v_n=-w_n\pmod3\) ならば、\(w\) と \(v\) は反対周期の対である。両者は異なる二つのサイクルに属してよい。
\subsection{半周期反周期性}
\(T=2h\) かつすべての \(n\) で \(w_{n+h}=-w_n\pmod3\) のとき、\(w\) は\textbf{半周期反周期的}である。このとき \(w=H\concat(-H)\)、\(H=(w_0,\ldots,w_{h-1})\)。例は \(01120221=0112\concat0221\)、\(011022=011\concat022\)、\(01220211=0122\concat0211\) である。
\subsection{後半の反転}
\(\mathcal R(a_0,\ldots,a_{h-1})=(a_{h-1},\ldots,a_0)\) は符号反転とは独立である。フィボナッチの例では \(\mathcal R(0221)=1220\)。したがって \(1220\) は \(0112\) の\textbf{反転された対蹠語}であり、後半そのものではない。
\section{唯一の反対シフト定理}
\begin{summarybox}\textbf{定理.} \(w\) を周期 \(T\) の非零原始語とする。すべての \(n\) で \(w_{n+r}=-w_n\) となるシフト \(r\) が存在するなら、\(T\) は偶数で \(r\equiv T/2\pmod T\) である。\end{summarybox}
\subsection*{証明}
シフトを二回適用すると \(w_{n+2r}=w_n\) となる。原始性より \(T\mid2r\)。非零の \(\F_3\) 語は \(w=-w\) を満たさないので \(r\not\equiv0\pmod T\)。巡回群 \(\Z/T\Z\) の非自明な位数 2 の元は \(T\) が偶数のときにのみ存在し、\(T/2\) に等しい。\hfill\(\square\)
\subsection*{帰結}
したがって非零原始周期は、\begin{enumerate}\item 符号反転によって別のサイクルへ送られる、または\item 符号反転で不変であり、その場合は必ず半周期反周期的である。\end{enumerate} この二分法は \(00111201,00222102\) のような対と \(011022\) のような内在的反周期を区別する。
\section{組合せ論的帰結}
\(w=H\concat(-H)\) なら、1 と 2 の個数は等しく、0 の個数は偶数で、\(\sum_{n=0}^{T-1}w_n=0\) が \(\F_3\) で成り立つ。これらは必要条件だが十分条件ではない。中心持ち上げ \(\widetilde0=0\)、\(\widetilde1=1\)、\(\widetilde2=-1\) により、生成多項式は \[W(X)=\sum_{n=0}^{2h-1}\widetilde w_nX^n=(1-X^h)\sum_{n=0}^{h-1}\widetilde w_nX^n\] と因数分解される。
\section{フーリエ特徴付け}
\(\widehat w(k)=\sum_{n=0}^{2h-1}\widetilde w_ne^{-2\pi i kn/(2h)}\) とおく。すると \(w_{n+h}=-w_n\) であることと、すべての偶数 \(k\) で \(\widehat w(k)=0\) であることは同値である。順方向では因子 \(1-(-1)^k\) が現れ、偶数 \(k\) で消える。逆に偶数モードがすべて消えれば語は奇数モードの部分空間に属し、\(h\) シフトは \(-I\) として作用する。これは目視比較ではなく厳密なスペクトル判定である。
\section{法 3 の線形数列}
\(k\) 次漸化式 \[u_{n+k}=c_0u_n+\cdots+c_{k-1}u_{n+k-1}\pmod3\] を考える。状態は \(s_n=(u_n,\ldots,u_{n+k-1})^{\mathsf T}\)、発展は \(s_{n+1}=Ms_n\) で、\(M\) はコンパニオン行列である。
\subsection{純周期性}
\(M\) の行列式は符号を除いて \(c_0\) である。\(\F_3\) 上で \(M\) が可逆であることと \(c_0\neq0\) は同値である。このときすべての状態軌道は純周期的である。\(c_0=0\) なら前周期が現れ得るため、サイクルと区別しなければならない。
\subsection{サイクルの符号反転}
線形性より \(M(-s)=-M(s)\)。よって符号反転は各サイクルを同じ長さのサイクルへ写す。非零サイクルでは、異なる二サイクルが交換されるか、一つのサイクルが不変で半周期反周期的になるかの二通りしかない。零サイクルは唯一の退化例外で、各点が固定され、非自明な原始反周期を与えない。
\subsection{一つの初期値に対する判定}
初期状態 \(s_0\) に対し \(M^hs_0=-s_0\) なら、任意の線形出力で \(u_{n+h}=-u_n\) が成り立つ。\(s_0\neq0\) かつサイクルが原始的なら、これは上記の半周期反周期を与える。
\subsection{一様判定}
すべての初期状態が同じ反周期関係 \(h\) を満たすことと \(M^h=-I\) は同値である。さらに最小多項式は \(\mu_M(X)\mid X^h+1\) を \(\F_3[X]\) で満たし、\(M^{2h}=I\) が従う。
\section{フィボナッチの例}
\(M=\begin{pmatrix}0&1\\1&1\end{pmatrix}\) を \(\F_3\) 上で考えると \(M^4=-I\)、\(M^8=I\)。初期値 \((0,1)\) は \(01120221=0112\concat(-0112)\) を生成する。後半の反転は \(\mathcal R(0221)=1220\)。十進表記での \(1220=L_{14}+F_{14}=2F_{15}\) は付加的な符号化恒等式であり、反周期性の一般的帰結ではない。
\section{任意の法 3 数列への一般化}
与えられた周期語の解析に漸化式の仮定は不要である。任意の原始周期について、回転サイクル、加法逆元、アフィン対称性 \(w_{n+r}=aw_n+b\)、反転対称性 \(w_{r-n}=aw_n+b\)、反周期性、偶奇フーリエモードを計算できる。行列理論が必要なのは線形漸化式が明示されている場合だけである。
\section{一般の法への拡張}
\(\Z/m\Z\) 上でも自然な対合は \(x\mapsto-x\) であり、\(w_{n+h}=-w_n\pmod m\) はそのまま有効である。\(m=2\) では符号反転が恒等写像となり、反周期性は通常の周期性に退化する。\(m>2\) では区別は非自明である。「\(m\) の補数」は整数代表に依存するので代数的定義には適さず、不変な用語は「法 \(m\) の加法逆元」である。
\section{アルゴリズム監査}
\(k\) 次の法 \(m\) 漸化式の周期は有限状態空間 \((\Z/m\Z)^k\) 上で検出すべきであり、任意長の接頭辞で反復ブロックを探すべきではない。スクリプトは基数 \(m\) の圧縮状態符号化、単一初期値に対する Brent 法（時間 \(O(\mu+\lambda)\)、空間 \(O(1)\)）、全初期値の正確な関数グラフ分解（時間 \(O(m^k)\)）、係数族間の並列化、反対サイクルと反周期サイクルの別分類、\(M^h=-I\) の行列検査を用いる。Mathematica も外部ラッパーも不要である。
\section{再現可能な走査結果}
十五の名前付き族の走査では、tritetranacci 族について長さ 24 と 8 の異なる反対サイクル対と、反周期 \(011022\)、\(01220211\)、\(12\) が得られる。さらに 2 次から 6 次までの非零二値係数ベクトル \(119\) 個を走査する：\(\sum_{k=2}^{6}(2^k-1)=119\)。結果は、\(M^h=-I\) を満たす族が 13、非自明な半周期反周期サイクルを少なくとも一つ持つ族が 88、異なる反対サイクル対を少なくとも一つ持つ族が 96 である。JSON と CSV から再現できる。
\section*{結論}\addcontentsline{toc}{section}{結論}
法 3 の符号反転は、周期語と線形漸化式のサイクル上の構造的対合として作用する。非零原始語がこの対合で不変になるのは、唯一の半周期シフトによる場合に限られる。同じ性質は生成多項式の因数分解、偶数フーリエモードの消失、または線形の場合の行列関係 \(M^h=-I\) によって同値に認識できる。したがってアルゴリズム解析は、サイクル固有の対称性と反対サイクル間の単なる対応を明確に分離する。
\end{document}